% ============================================================================
\section{Evolution of the Stacked Polyphase Bridge Concept}
\label{sec:evolution}
% ============================================================================

This section traces how SPB's converter--machine modularity, together
with its associated control and modulation methods, has developed over
more than a decade of research, from an initial concept to an
increasingly reconfigurable and actively controlled
converter--machine system.

SPB emerged from a broader development toward modular converter--machine
systems. Earlier modular and integrated modular motor drives (IMMDs)
established the building blocks SPB later drew on, repeated
converter-winding modules, machine segmentation, power sharing, and
fault-tolerant operation \cite{Cengelci2000MMMI,Choi2011IMMD}, and
series-connected IMMD modules further showed that the total dc voltage
could be shared among inverter segments to enable lower-voltage
devices \cite{Wang2013GaNIMMD,Wang2025GaNIMMD}. A 2015 IEEE TIA study
by Wang, Li, and Han is especially relevant to SPB's later work,
combining series-connected GaN IMMD hardware with active cell-voltage
balancing and interleaved switching \cite{Wang2015GaNIMMD}.

Independent of this line of work, other series-connected modular
motor-drive topologies were proposed with the same underlying
dc-voltage-sharing principle \cite{Wang2015SeriesModular,
Ruseler2015InputSeries}, indicating that the concept was not unique
to SPB but part of a broader convergence toward series-connected
modularity. A related KTH comparative study evaluated SPB alongside
other integrated modular converter candidates for traction
drives, the parallel-connected polyphase bridge converter (PPB-I, a
conventional multiphase converter with phase-leg-level modularity, and
PPB-II, several interleaved two-level three-phase converters each with
its own dc capacitor) and the modular high-frequency (MHF) converter,
against a conventional two-level three-phase drive as a baseline in
terms of machine design, losses, capacitor energy storage, cost, and
redundancy, confirming that SPB is one of several modular-converter
topology candidates studied for this application
\cite{Zhang2017ModularDriveComparison}.

Building on this foundation, the SPB topology itself was introduced in
2013 \cite{Norrga2013SPB}, followed shortly thereafter by dedicated
control and modulation work \cite{Jin2014Control}. Related series-connected multiphase drives
demonstrated that machine-side degrees of freedom could be used to
regulate internal dc-link voltage sharing \cite{Che2014SeriesConnected},
while Nikouie \textit{et al.} formally established the SPB dc-link
stability problem and the corresponding balancing-controller
requirements \cite{Nikouie2015Stability,Nikouie2017Stability}. This
period established the topology's basic feasibility and identified
dc-link voltage balancing, rather than semiconductor voltage rating
alone, as its first-order control challenge.

The next stage of development shifted the focus toward modulation,
scalable control, and converter--machine co-design: modulation was
shown to directly influence semiconductor loss and dc-link harmonic
content \cite{Jin2017Modulation}, distributed control offered an
alternative to centralized balancing \cite{Jin2018Distributed}, and
interleaving and phase-count studies tied converter design to machine
architecture \cite{Verkroost2019Interleaving,VanDamme2019DCLink,
Bringezu2020SPB,Verkroost2021Ripple}. Collectively, this reframed SPB
from a series-connected converter topology into a converter--machine
co-design problem in which modulation, control, and machine parameters
jointly determine dc-link behavior.

More recent work has moved from repeated but statically configured
hardware toward actively reconfigurable modular systems.
Common-duty-ratio operation, decentralized voltage balancing, and
multi-agent power distribution have reduced the dependence on
centralized control \cite{Mocevic2023CDR,Verkroost2024VoltageBalance,
Verkroost2024PowerDistribution,Mocevic2024GateDelay}. Online module
isolation and reintroduction have further demonstrated dynamic dc-bus
reconfiguration, removing a faulted cell without interrupting torque
production \cite{Verkroost2024Reconfiguration}.
In parallel, modern GaN devices have renewed interest in SPB as a
route to lower semiconductor voltage classes in integrated drives
\cite{Rohner2024GaN}.

Fig.~\ref{fig:spb_timeline} places representative milestones from
across these stages on a chronological timeline of the topology's
development, together with the validation level, reported scale, and
the specific implication each result carries for the system-level
review developed in this paper.

\begin{figure}[!t]
\centering
\begin{tikzpicture}[font=\scriptsize,>=latex]
\tikzstyle{yr}=[circle,draw=blue!65!black,fill=blue!55!black,inner sep=1.2pt]
\tikzstyle{txt}=[align=left,text width=48mm,inner sep=0pt]
\draw[->,thick,blue!55!black] (0,3mm) -- (0,-53mm);
\node[yr] (n1) at (0,0mm) {};
\node[txt,right=2mm of n1]
{\textbf{2013} -- Norrga \textit{et al.}~\cite{Norrga2013SPB}:
original SPB topology introduced};
\node[yr] (n2) at (0,-10mm) {};
\node[txt,right=2mm of n2]
{\textbf{2017} -- Jin \textit{et al.}~\cite{Jin2017Modulation}:
modulation shown to affect switching loss and dc-link harmonics};
\node[yr] (n3) at (0,-20mm) {};
\node[txt,right=2mm of n3]
{\textbf{2023} -- Mocevic \textit{et al.}~\cite{Mocevic2023CDR}:
CDR balancing demonstrated for an induction-machine system};
\node[yr] (n4) at (0,-30mm) {};
\node[txt,right=2mm of n4]
{\textbf{2024} -- Verkroost \textit{et al.}
~\cite{Verkroost2024VoltageBalance,Verkroost2024PowerDistribution}:
distributed balancing and power sharing, 4~kW axial-flux PMSM};
\node[yr] (n5) at (0,-40mm) {};
\node[txt,right=2mm of n5]
{\textbf{2024} -- Verkroost \textit{et al.}
~\cite{Verkroost2024Reconfiguration}: online module
isolation/reactivation, $\sim$100~ms reconfiguration};
\node[yr] (n6) at (0,-50mm) {};
\node[txt,right=2mm of n6]
{\textbf{2024} -- Rohner \textit{et al.}~\cite{Rohner2024GaN}:
650-V GaN SPB vs.\ 3L flying-capacitor comparison, 7.5~kW, 800~V};
\end{tikzpicture}
\caption{Chronological milestones in SPB research, from the original
topology through modulation, distributed balancing, reconfiguration,
and wide-bandgap integration.}
\label{fig:spb_timeline}
\end{figure}

As Fig.~\ref{fig:spb_timeline} indicates, this progression has been
validated primarily at laboratory scale, typically at or below several
kilowatts, leaving open whether the demonstrated control, balancing,
and reconfiguration principles extend to the automotive and heavy-duty
power levels summarized in Table~\ref{tab:intro_voltage_power} -- a
gap revisited throughout the remainder of this paper.